\documentclass[12pt]{article}
\usepackage[T1]{fontenc}
\usepackage[utf8]{inputenc}
\usepackage{lmodern}
\usepackage{textcomp}
\usepackage{enumitem}
\usepackage{geometry}
\usepackage{graphicx}
\usepackage{amsmath, amssymb}
\geometry{margin=1in}
\begin{document}
\begin{center}
Chemistry - Graduate Level Assessment 14 \
Total Marks: 40 \
Answer all questions.
\end{center}

\section*{Multiple Choice (10 marks)}
\begin{enumerate}[label=\arabic*.]
\item An automobile travelling at 10 miles per hour produces 0.33 1 b of CO gas per mile. How many moles of CO are produced per mile?
\begin{enumerate}[label=(\alph*)]
    \item 5.4 moles per mile
    \item 9.0 moles per mile
    \item 2.5 moles per mile
    \item 7.3 moles per mile
    \item 3.8 moles per mile
\end{enumerate}
\item What weight of sulfur must combine with aluminum to form 600 lbs of aluminum sulfide?
\begin{enumerate}[label=(\alph*)]
    \item 350 lbs. sulfur
    \item 320 lbs. sulfur
    \item 425 lbs. sulfur
    \item 400 lbs. sulfur
    \item 300 lbs. sulfur
\end{enumerate}
\item Sirius, one of the hottest known stars, has approximately a blackbody spectrum with $\lambda_{\max }=260 \mathrm{~nm}$. Estimate the surface temperature of Sirius.
\begin{enumerate}[label=(\alph*)]
    \item 8500 $\mathrm{~K}$
    \item 15000 $\mathrm{~K}$
    \item 9500 $\mathrm{~K}$
    \item 11000 $\mathrm{~K}$
\end{enumerate}
\item SO\_2 CI\_2 \ensuremath{\rightarrow} SO\_2 + CI\_2 is 8.0 minutes. In what period of time would the con- centration of SO\_2 CI\_2 be reduced to 1.0\% of the original?
\begin{enumerate}[label=(\alph*)]
    \item 43 minutes
    \item 103 minutes
    \item 93 minutes
    \item 113 minutes
    \item 23 minutes
\end{enumerate}
\item A gas of molecular weight 58 was found to contain 82.8\% C and 17.2\% H. What is the molecular formula of this com-pound?
\begin{enumerate}[label=(\alph*)]
    \item C\_3 H\_6
    \item C\_5 H\_10
    \item C\_4 H\_8
    \item C\_4 H\_9
    \item C\_3 H\_8
\end{enumerate}
\end{enumerate}

\section*{True / False (10 marks)}
\textit{Answer True or False}
\begin{enumerate}[label=\arabic*.]
\setcounter{enumi}{5}
\item True or False: The final concentration of the mixed solution is 0.45 M H$_{2}$SO$_{4}$, indicating an increase in molarity due to dilution.
\item True or False: The number of configuration state functions in a full CI calculation of CH 2 SiHF using a 6-31 G** basis set is 9.45 x $10^27$.
\item True or False: Reacting 80 g of O\_2 requires 72 g of NH\_3, as it is a common misconception to overestimate the amount needed.
\item True or False: If the position of an electron is known within a 10 pm interval, the uncertainty in its momentum is 9.1 x 10^\{-23\} kg·m·s^\{-1\}.
\item True or False: During the boiling of liquid NH 3, the hydrogen bonds holding separate NH 3 molecules together break apart.
\end{enumerate}

\section*{Long Form Response (20 marks)}
\textit{Answer each question with a clear and complete explanation.}
\begin{enumerate}[label=\arabic*.]
\setcounter{enumi}{10}
\item Discuss the relationship between kinetic energy and wavelength for a photon and an electron, each possessing 1 eV of kinetic energy. Include calculations and implications of these wavelengths.
\item Discuss the weight of 1,000 cubic feet of air at Standard Temperature and Pressure (STP), including the calculations and implications of your findings in a broader scientific context.
\end{enumerate}

\vfill
\noindent\textit{End of Paper}
\end{document}